\newpage
\section{综合习题}
以下习题涵盖本章的主要内容，不再按小节分类。解题时应注意定义域与各公式的适用条件，并说明所用的函数性质。

\begin{exercise}
\label{exs:functions-comprehensive}
\mbox{}\par\nobreak

（1）求下列函数的定义域：

\quad \exercisesubitem{1}\(f(x) = \sqrt{\ln \left( 2 x - 1 \right)}\)；

\quad \exercisesubitem{2}\(g(x) = \ln \frac{x + 1}{x - 1}\)；

\quad \exercisesubitem{3}\(h(x) = \sqrt{1 - 2 \sin x}\)。

（2）设\(f(x) = x^{2}\)，定义域为\(A = [-2 , 1]\)。

\quad \exercisesubitem{1}求\(f\)的值域，判断其是否有界，是否为奇函数或偶函数，以及是否存在反函数。

\quad \exercisesubitem{2}将定义域限制为\([-2 , 0]\)，求此时的反函数及其定义域、值域。说明同一个表达式在不同定义域下，函数的性质会发生哪些变化。

（3）设\(f(x) = \ln x\)，\(g(x) = x^{2} - 1\)，两者均取使表达式有意义的最大实数定义域。

\quad \exercisesubitem{1}求\(f \circ g\)与\(g \circ f\)的表达式、定义域和值域。

\quad \exercisesubitem{2}判断这两个复合函数是否相同，并求它们的单调区间。

（4）设\(f(x) = a x + b\)，\(g(x) = c x + d\)，定义域均为\(\mathbb{R}\)。

\quad \exercisesubitem{1}求\(f \circ g = g \circ f\)成立的充要条件。

\quad \exercisesubitem{2}举出\(f \ne g\)但复合运算可交换的例子，再举出不可交换的例子。

\quad \exercisesubitem{3}当\(a \ne 0\)时求\(f^{-1}\)，并验证\(f \circ f^{-1}\)与\(f^{-1} \circ f\)的作用。

（5）考虑方程\(\left( x - 1 \right)^{2} + y^{2} = 4\)。

\quad \exercisesubitem{1}这个方程是否在\([-1 , 3]\)上确定了一个以\(x\)为自变量的函数？若再要求\(y \ge 0\)，写出对应函数的显式表示、定义域和值域。

\quad \exercisesubitem{2}曲线还可以用参数方程
\begin{equation*}
    \left\{ \begin{aligned}
        x &= 1 + 2 \cos t , \\
        y &= 2 \sin t
    \end{aligned} \right.
\end{equation*}
表示。分别取\(t \in [0 , \pi]\)与\(t \in [0 , 2 \pi]\)，说明所表示的图形，以及它们能否整体看成以\(x\)为自变量的函数图像。

（6）设
\begin{equation*}
    f(x) = \begin{cases}
        x^{2} , & x < 1 , \\
        a x + b , & x \ge 1 .
    \end{cases}
\end{equation*}

\quad \exercisesubitem{1}求\(f\)在\(x = 1\)处连续的充要条件。

\quad \exercisesubitem{2}取\(a = 2\)，\(b = -1\)，画出函数图像，求其值域和最大单调区间，并判断其在\(\mathbb{R}\)上是否有界、是否存在反函数。

（7）分别讨论\(x^{-2}\)、\(x^{1/3}\)、\(x^{\sqrt{2}}\)与\(x^{0}\)在基本定义域\(x > 0\)上的连续性、单调性和值域。

再讨论这些幂函数能否扩展到\(x < 0\)或\(x = 0\)，以及扩展后能否讨论奇偶性。对于无理数次幂，说明它与有理数次幂之间的联系。

（8）利用双曲函数的定义完成下列问题：

\quad \exercisesubitem{1}判断\(\sinh x\)、\(\cosh x\)与\(\tanh x\)的奇偶性，并证明\(\cosh^{2} x - \sinh^{2} x = 1\)。

\quad \exercisesubitem{2}求\(\sinh x\)、\(\cosh x\)与\(\tanh x\)的值域，并判断哪些函数在\(\mathbb{R}\)上有界。

\quad \exercisesubitem{3}证明\(\sinh x\)严格递增，求其反函数；\(\cosh x\)在\(\mathbb{R}\)上是否有反函数？若限制在\([0 , + \infty)\)上，情形如何？

（9）设\(A\)是一个关于原点对称的非空实数集合，\(f\)是\(A\)上的任意实值函数。定义
\begin{equation*}
    u(x) = \frac{f(x) + f(-x)}{2} , \qquad
    v(x) = \frac{f(x) - f(-x)}{2} .
\end{equation*}

\quad \exercisesubitem{1}证明\(u\)为偶函数，\(v\)为奇函数，且\(f = u + v\)。

\quad \exercisesubitem{2}这种“一个偶函数与一个奇函数之和”的表示是否唯一？

\quad \exercisesubitem{3}若\(f\)在\(A\)上有界，\(u\)与\(v\)是否也有界？若\(f\)连续，它们是否也连续？

（10）列出\(\tan x\)、\(\cot x\)、\(\sec x\)与\(\csc x\)的定义域、值域、奇偶性与最小正周期。

分别求这四个函数在\((0 , \pi)\)内的最大单调区间，注意将没有定义的点排除。

（11）设\(f(x) = 2 \sin \left( 3 x - \frac{\pi}{4} \right) + 1\)。

\quad \exercisesubitem{1}说明可以怎样从\(y = \sin x\)的图像得到\(f\)的图像，并求\(f\)的值域与最小正周期。

\quad \exercisesubitem{2}求\(f\)的最大单调增区间、最大单调减区间，以及图像的全部对称轴和对称中心。

\quad \exercisesubitem{3}判断\(f\)是否为奇函数或偶函数，并求它在\([0 , \pi]\)上取最大值的所有自变量。

（12）设
\begin{equation*}
    f(x) = \ln \left( 3 + \sin 2 x + \sqrt{3} \cos 2 x \right) .
\end{equation*}

\quad \exercisesubitem{1}利用和差角公式化简对数内部的表达式，求\(f\)的定义域、值域与最小正周期，并说明其连续性。

\quad \exercisesubitem{2}将\(f\)限制在区间\(I = \left[ - \frac{5 \pi}{12} , \frac{\pi}{12} \right]\)上，说明为什么此时存在反函数，并求反函数及其定义域、值域。

\quad \exercisesubitem{3}求方程\(f(x) = \ln 4\)的全部实数解，并指出其中位于\(I\)内的解。

（13）利用三角公式完成下列问题：

\quad \exercisesubitem{1}推导\(\sin 3 x = 3 \sin x - 4 \sin^{3} x\)。

\quad \exercisesubitem{2}证明
\begin{equation*}
    \frac{\sin \alpha + \sin \beta}{\cos \alpha + \cos \beta}
    = \tan \frac{\alpha + \beta}{2} ,
\end{equation*}
并说明等式的适用条件。

\quad \exercisesubitem{3}在\([0 , 2 \pi]\)内求解方程\(\cos 2 x = \sin x\)，以及不等式\(\sin x + \cos x \ge 1\)。

（14）在\(\triangle A B C\)中，\(b = 2\)，\(c = 3\)，\(A = \frac{\pi}{3}\)。

\quad \exercisesubitem{1}求边长\(a\)、三角形的面积及外接圆半径。

\quad \exercisesubitem{2}保持\(b = 2\)、\(c = 3\)不变，让\(A\)在\((0 , \pi)\)内变化。讨论边长\(a\)如何随\(A\)变化，并求面积的最大值及取到最大值时的\(A\)。

（15）在\(\triangle A B C\)中，\(a = 1\)，\(b = \sqrt{2}\)，\(A = \frac{\pi}{6}\)。

\quad \exercisesubitem{1}这样的三角形有几个？分别求所有可能的\(B\)、\(C\)与\(c\)。

\quad \exercisesubitem{2}说明为什么由正弦定理求出\(\sin B\)之后，不能直接认定\(B\)只有一个取值，并与已知两边及其夹角的情形比较。
\end{exercise}
